%
% workspace-model.tex
%
% Z Specification for the Lux Workspace Model (DES-022)
% Formal model of frames, client identity, world menu, and frame aggregation
%

\documentclass[a4paper,10pt,fleqn]{article}
\usepackage[margin=1in]{geometry}
\usepackage{fuzz}
\usepackage[colorlinks=true,linkcolor=blue,citecolor=blue,urlcolor=blue]{hyperref}

\begin{document}

\title{Lux Workspace Model: A Z Specification}
\author{Formal Model of Frames, Client Identity, and World Menu}
\date{March 2026}
\hypersetup{
  pdftitle={Lux Workspace Model: A Z Specification},
  pdfauthor={Punt Labs},
  pdfsubject={Z Specification},
  pdfcreator={fuzz/probcli}
}
\maketitle

\tableofcontents
\newpage

% ===================================================================
\section{Introduction}
% ===================================================================

DES-022 introduces a workspace model for the Lux display server,
inspired by Pharo/Squeak's live environment.  The display becomes a
\emph{workspace} containing independent, movable \emph{frames} (ImGui
inner windows).  Each frame is owned by a connected client and
contains one or more \emph{scenes}.

This specification models the workspace layer that sits above the
existing scene renderer (specified in \texttt{display-server.tex}).
It covers:
\begin{itemize}
  \item Client identity from the connection handshake
  \item Frame lifecycle: creation, minimization, restoration, closure
  \item Frame ownership and cleanup on client disconnect
  \item World menu: per-client namespaced menu registration
  \item Scene-to-frame targeting
\end{itemize}

We abstract away the element tree inside scenes (already modeled) and
focus on the new state that DES-022 introduces.

% ===================================================================
\section{Basic Types}
% ===================================================================

\begin{zed}
  [FD, FRAMEID, SCENEID, CLIENTNAME, MENUID, MENUPATH]
\end{zed}

$FD$ models socket file descriptors (shared with the display-server
spec).  $FRAMEID$ is the identifier a client assigns to a frame.
$SCENEID$ identifies a scene within a frame.  $CLIENTNAME$ is the
human-readable display name declared during the connection handshake
(e.g.\ ``Lux'', ``Vox'', ``Quarry'').  $MENUID$ and $MENUPATH$
identify menu items and their nesting path in the World menu.

% ===================================================================
\section{Free Types}
% ===================================================================

\begin{zed}
  ZBOOL ::= ztrue | zfalse
\end{zed}

Frame visibility state:
\begin{zed}
  FrameVis ::= fvNormal | fvMinimized
\end{zed}

Menu target --- which menu surface an item belongs to:
\begin{zed}
  MenuTarget ::= mtTools | mtWorld
\end{zed}

% ===================================================================
\section{Global Constants}
% ===================================================================

Animation-friendly bounds.  Production values would be larger;
these are sized for ProB's finite state space.

\begin{axdef}
  maxClients : \nat \\
  maxFrames : \nat \\
  maxScenes : \nat \\
  maxMenuItems : \nat
  \where
  maxClients = 3 \\
  maxFrames = 4 \\
  maxScenes = 6 \\
  maxMenuItems = 6
\end{axdef}

% ===================================================================
\section{Workspace State}
% ===================================================================

The workspace tracks connected clients with their identities,
frames with their ownership and visibility, scene-to-frame
mappings, and per-client menu registrations.  All fields are
flattened for ProB animatability.

\begin{schema}{Workspace}
  clients : \power FD \\
  clientNames : FD \pfun CLIENTNAME \\
  frames : \power FRAMEID \\
  frameOwner : FRAMEID \pfun FD \\
  frameVis : FRAMEID \pfun FrameVis \\
  frameScenes : FRAMEID \pfun \power SCENEID \\
  sceneFrame : SCENEID \pfun FRAMEID \\
  menuItems : \power MENUID \\
  menuOwner : MENUID \pfun FD \\
  menuTarget : MENUID \pfun MenuTarget \\
  menuPath : MENUID \pfun MENUPATH
  \where
  % --- Client identity ---
  \dom clientNames \subseteq clients \\
  \# clients \leq maxClients \\
  % --- Frame ownership ---
  \dom frameOwner = frames \\
  \dom frameVis = frames \\
  \dom frameScenes = frames \\
  \ran frameOwner \subseteq clients \\
  \# frames \leq maxFrames \\
  % --- Scene-frame bijection ---
  \dom sceneFrame = \bigcup (\ran frameScenes) \\
  \forall f : frames @ \forall s : frameScenes~f @
  \quad~ sceneFrame~s = f \\
  \#(\dom sceneFrame) \leq maxScenes \\
  % --- Menu ownership ---
  \dom menuOwner = menuItems \\
  \dom menuTarget = menuItems \\
  \dom menuPath = menuItems \\
  \ran menuOwner \subseteq clients \\
  \# menuItems \leq maxMenuItems
\end{schema}

The key invariants are:
\begin{enumerate}
  \item \textbf{Named clients are connected}: Only connected
    clients can have display names.
  \item \textbf{Frame owners are connected}: Every frame is owned
    by a connected client.  Disconnecting a client removes its frames.
  \item \textbf{Frame metadata coverage}: Every frame has an owner,
    visibility state, and scene set.
  \item \textbf{Scene-frame consistency}: Every scene knows which
    frame it belongs to, and that mapping is consistent with the
    frame's scene set.
  \item \textbf{Menu owners are connected}: Every menu item is
    owned by a connected client.
\end{enumerate}

% ===================================================================
\section{Initialization}
% ===================================================================

\begin{schema}{Init}
  Workspace'
  \where
  clients' = \emptyset \\
  clientNames' = \emptyset \\
  frames' = \emptyset \\
  frameOwner' = \emptyset \\
  frameVis' = \emptyset \\
  frameScenes' = \emptyset \\
  sceneFrame' = \emptyset \\
  menuItems' = \emptyset \\
  menuOwner' = \emptyset \\
  menuTarget' = \emptyset \\
  menuPath' = \emptyset
\end{schema}

% ===================================================================
\section{Operations}
% ===================================================================

% -------------------------------------------------------------------
\subsection{Accept Connection}
% -------------------------------------------------------------------

A new client connects.  It has no display name yet (the
\texttt{ConnectMessage} arrives later).

\begin{schema}{AcceptConnection}
  \Delta Workspace \\
  newClient? : FD
  \where
  newClient? \notin clients \\
  \# clients < maxClients \\
  clients' = clients \cup \{ newClient? \} \\
  clientNames' = clientNames \\
  frames' = frames \\
  frameOwner' = frameOwner \\
  frameVis' = frameVis \\
  frameScenes' = frameScenes \\
  sceneFrame' = sceneFrame \\
  menuItems' = menuItems \\
  menuOwner' = menuOwner \\
  menuTarget' = menuTarget \\
  menuPath' = menuPath
\end{schema}

% -------------------------------------------------------------------
\subsection{Identify Client}
% -------------------------------------------------------------------

A connected client declares its display name via
\texttt{ConnectMessage}.  Subsequent calls update the name (idempotent).

\begin{schema}{IdentifyClient}
  \Delta Workspace \\
  fd? : FD \\
  name? : CLIENTNAME
  \where
  fd? \in clients \\
  clientNames' = clientNames \oplus \{ fd? \mapsto name? \} \\
  clients' = clients \\
  frames' = frames \\
  frameOwner' = frameOwner \\
  frameVis' = frameVis \\
  frameScenes' = frameScenes \\
  sceneFrame' = sceneFrame \\
  menuItems' = menuItems \\
  menuOwner' = menuOwner \\
  menuTarget' = menuTarget \\
  menuPath' = menuPath
\end{schema}

% -------------------------------------------------------------------
\subsection{Create Frame}
% -------------------------------------------------------------------

A client sends a \texttt{SceneMessage} with a \texttt{frame\_id} that
does not yet exist.  The display creates a new frame containing the
specified scene.

\begin{schema}{CreateFrame}
  \Delta Workspace \\
  fd? : FD \\
  frameId? : FRAMEID \\
  sceneId? : SCENEID
  \where
  fd? \in clients \\
  frameId? \notin frames \\
  sceneId? \notin \dom sceneFrame \\
  \# frames < maxFrames \\
  \#(\dom sceneFrame) < maxScenes \\
  frames' = frames \cup \{ frameId? \} \\
  frameOwner' = frameOwner \cup \{ frameId? \mapsto fd? \} \\
  frameVis' = frameVis \cup \{ frameId? \mapsto fvNormal \} \\
  frameScenes' = frameScenes \cup
  \quad~ \{ frameId? \mapsto \{ sceneId? \} \} \\
  sceneFrame' = sceneFrame \cup \{ sceneId? \mapsto frameId? \} \\
  clients' = clients \\
  clientNames' = clientNames \\
  menuItems' = menuItems \\
  menuOwner' = menuOwner \\
  menuTarget' = menuTarget \\
  menuPath' = menuPath
\end{schema}

% -------------------------------------------------------------------
\subsection{Add Scene to Frame}
% -------------------------------------------------------------------

A client sends a \texttt{SceneMessage} targeting an existing frame
with a new scene.  The scene is added (creating a tab if this is
the second scene in the frame).

\begin{schema}{AddSceneToFrame}
  \Delta Workspace \\
  fd? : FD \\
  frameId? : FRAMEID \\
  sceneId? : SCENEID
  \where
  fd? \in clients \\
  frameId? \in frames \\
  frameOwner~frameId? = fd? \\
  sceneId? \notin \dom sceneFrame \\
  \#(\dom sceneFrame) < maxScenes \\
  frameScenes' = frameScenes \oplus
  \quad~ \{ frameId? \mapsto (frameScenes~frameId? \cup \{ sceneId? \}) \} \\
  sceneFrame' = sceneFrame \cup \{ sceneId? \mapsto frameId? \} \\
  frames' = frames \\
  frameOwner' = frameOwner \\
  frameVis' = frameVis \\
  clients' = clients \\
  clientNames' = clientNames \\
  menuItems' = menuItems \\
  menuOwner' = menuOwner \\
  menuTarget' = menuTarget \\
  menuPath' = menuPath
\end{schema}

% -------------------------------------------------------------------
\subsection{Update Scene in Frame}
% -------------------------------------------------------------------

A client sends a \texttt{SceneMessage} targeting an existing frame
with an existing scene ID.  The scene content is replaced (modeled
  here as a no-op on workspace state since scene \emph{content}
is in the display-server spec).

\begin{schema}{UpdateScene}
  \Xi Workspace \\
  fd? : FD \\
  frameId? : FRAMEID \\
  sceneId? : SCENEID
  \where
  fd? \in clients \\
  frameId? \in frames \\
  frameOwner~frameId? = fd? \\
  sceneId? \in frameScenes~frameId?
\end{schema}

% -------------------------------------------------------------------
\subsection{Close Frame (User)}
% -------------------------------------------------------------------

The user clicks the close button on a frame.  The frame and all
its scenes are removed.  An interaction event is sent to the
owning client (modeled at the display-server level).

\begin{schema}{CloseFrame}
  \Delta Workspace \\
  frameId? : FRAMEID
  \where
  frameId? \in frames \\
  frames' = frames \setminus \{ frameId? \} \\
  frameOwner' = \{ frameId? \} \ndres frameOwner \\
  frameVis' = \{ frameId? \} \ndres frameVis \\
  frameScenes' = \{ frameId? \} \ndres frameScenes \\
  sceneFrame' = (frameScenes~frameId?) \ndres sceneFrame \\
  clients' = clients \\
  clientNames' = clientNames \\
  menuItems' = menuItems \\
  menuOwner' = menuOwner \\
  menuTarget' = menuTarget \\
  menuPath' = menuPath
\end{schema}

% -------------------------------------------------------------------
\subsection{Minimize Frame}
% -------------------------------------------------------------------

\begin{schema}{MinimizeFrame}
  \Delta Workspace \\
  frameId? : FRAMEID
  \where
  frameId? \in frames \\
  frameVis~frameId? = fvNormal \\
  frameVis' = frameVis \oplus \{ frameId? \mapsto fvMinimized \} \\
  frames' = frames \\
  frameOwner' = frameOwner \\
  frameScenes' = frameScenes \\
  sceneFrame' = sceneFrame \\
  clients' = clients \\
  clientNames' = clientNames \\
  menuItems' = menuItems \\
  menuOwner' = menuOwner \\
  menuTarget' = menuTarget \\
  menuPath' = menuPath
\end{schema}

% -------------------------------------------------------------------
\subsection{Restore Frame}
% -------------------------------------------------------------------

\begin{schema}{RestoreFrame}
  \Delta Workspace \\
  frameId? : FRAMEID
  \where
  frameId? \in frames \\
  frameVis~frameId? = fvMinimized \\
  frameVis' = frameVis \oplus \{ frameId? \mapsto fvNormal \} \\
  frames' = frames \\
  frameOwner' = frameOwner \\
  frameScenes' = frameScenes \\
  sceneFrame' = sceneFrame \\
  clients' = clients \\
  clientNames' = clientNames \\
  menuItems' = menuItems \\
  menuOwner' = menuOwner \\
  menuTarget' = menuTarget \\
  menuPath' = menuPath
\end{schema}

% -------------------------------------------------------------------
\subsection{Register Menu Item}
% -------------------------------------------------------------------

A client registers a menu item.  The \texttt{menuTarget} determines
whether it appears in the Tools menu or the World menu
(per-client namespace from handshake identity).

\begin{schema}{RegisterMenuItem}
  \Delta Workspace \\
  fd? : FD \\
  itemId? : MENUID \\
  target? : MenuTarget \\
  path? : MENUPATH
  \where
  fd? \in clients \\
  itemId? \notin menuItems \\
  \# menuItems < maxMenuItems \\
  menuItems' = menuItems \cup \{ itemId? \} \\
  menuOwner' = menuOwner \cup \{ itemId? \mapsto fd? \} \\
  menuTarget' = menuTarget \cup \{ itemId? \mapsto target? \} \\
  menuPath' = menuPath \cup \{ itemId? \mapsto path? \} \\
  clients' = clients \\
  clientNames' = clientNames \\
  frames' = frames \\
  frameOwner' = frameOwner \\
  frameVis' = frameVis \\
  frameScenes' = frameScenes \\
  sceneFrame' = sceneFrame
\end{schema}

% -------------------------------------------------------------------
\subsection{Deregister Menu Item}
% -------------------------------------------------------------------

A client removes a menu item it owns.

\begin{schema}{DeregisterMenuItem}
  \Delta Workspace \\
  fd? : FD \\
  itemId? : MENUID
  \where
  fd? \in clients \\
  itemId? \in menuItems \\
  menuOwner~itemId? = fd? \\
  menuItems' = menuItems \setminus \{ itemId? \} \\
  menuOwner' = \{ itemId? \} \ndres menuOwner \\
  menuTarget' = \{ itemId? \} \ndres menuTarget \\
  menuPath' = \{ itemId? \} \ndres menuPath \\
  clients' = clients \\
  clientNames' = clientNames \\
  frames' = frames \\
  frameOwner' = frameOwner \\
  frameVis' = frameVis \\
  frameScenes' = frameScenes \\
  sceneFrame' = sceneFrame
\end{schema}

% -------------------------------------------------------------------
\subsection{Disconnect Client}
% -------------------------------------------------------------------

A client disconnects.  All its frames, scenes, and menu items
are removed.  This is the critical cleanup operation that
maintains the ownership invariants.

Let $deadFrames$ be the set of frames owned by the departing client,
and $deadScenes$ be all scenes in those frames, and $deadMenus$
be all menu items owned by the client.

\begin{schema}{DisconnectClient}
  \Delta Workspace \\
  deadClient? : FD
  \where
  deadClient? \in clients \\
  clients' = clients \setminus \{ deadClient? \} \\
  clientNames' = \{ deadClient? \} \ndres clientNames \\
  frames' = frames \setminus
  \quad~ \{ f : frames | frameOwner~f = deadClient? \} \\
  frameOwner' = frames' \dres frameOwner \\
  frameVis' = frames' \dres frameVis \\
  frameScenes' = frames' \dres frameScenes \\
  sceneFrame' = sceneFrame \rres frames' \\
  menuItems' = menuItems \setminus
  \quad~ \{ m : menuItems | menuOwner~m = deadClient? \} \\
  menuOwner' = menuItems' \dres menuOwner \\
  menuTarget' = menuItems' \dres menuTarget \\
  menuPath' = menuItems' \dres menuPath
\end{schema}

% ===================================================================
\section{System Invariants Summary}
% ===================================================================

The key properties maintained across all operations:

\begin{enumerate}
  \item \textbf{Frame owners are connected}: $\ran frameOwner
    \subseteq clients$.  No orphaned frames.

  \item \textbf{Scene-frame consistency}: Every scene maps to
    exactly one frame, and that frame lists the scene in its
    scene set.  No orphaned scenes.

  \item \textbf{Menu owners are connected}: $\ran menuOwner
    \subseteq clients$.  No orphaned menu items.

  \item \textbf{Disconnect atomicity}: When a client disconnects,
    all its frames, scenes, and menu items are removed in
    one atomic operation.

  \item \textbf{Bounded resources}: Client count, frame count,
    scene count, and menu item count are all bounded.

  \item \textbf{Named clients are connected}: $\dom clientNames
    \subseteq clients$.  Names cannot outlive connections.

  \item \textbf{Ownership exclusivity}: Each frame has exactly one
    owner ($frameOwner$ is a total function on $frames$).
    Only the owning client can add scenes to a frame.
\end{enumerate}

\end{document}
